\scanpage{129}
% archive sheet number: 128
pour $x\in E=\prod_{i\in I} X_i$, \deleted{et $g=(g_i)_{i\in I}\in V$}\inserted{et $i_0\in I$}
\[
\begin{aligned}
(\sigma_{i_0})_E\deleted{\unsure{(x)=}}&=\prod_{i\in I\smallsetminus\{i_0\}}\rho(\varepsilon_i(z_0))
=\rho\biggl(\underbrace{\sum_{i\in I\smallsetminus i_0}\varepsilon_i(z_0)}_{\unsure{z\ \varepsilon_{i_0}(z_0)|z+\varepsilon_{i_0}(z_0)}}\biggr)\\
&=\rho(z)\rho(\varepsilon_{i_0}(z_0))
\end{aligned}
\]

Donc si
\[
V\xrightarrow{\unsure{\eta|\gamma}}\mathbb F_2\qquad
(\text{on a aussi }\mathbb F_2\xrightarrow{\delta=\mathrm{diag}}V)
\]
est l'augmentation naturelle, on aura pour $g\in V$
\[
(\sigma_g)_E=\rho(g\deleted{\unsure{\rho}}+\delta\eta(g))
\]

Considérons l'application
\[
V\xrightarrow{u}V\qquad u(x)=x+\delta\eta(\deleted{\unsure{g}}x)
\]
Son noyau est formé des $x$ tels que
\[
x=\delta\eta x
\]
donc \deleted{\unsure{est contenu}}\inserted{l'unique} élément non nul possible de ce noyau est $z=\delta(z_0)$, et il reste à voir quand on a $\deleted{\unsure{\delta|z}}z=\delta\eta\deleted{\unsure{\delta}}z$ i.e. $\deleted{z_0=\eta(\delta z_0}$ (comme $z=\delta z_0$) $z_0=\eta z$ i.e. $\eta z\ne0$: c'est le cas \unsure{sss|ssi} $I$ est de cardinal impair. Or ce cardinal est $1+c_\nu=1+2^\nu$, qui est impair sauf pour $\nu=0$ (cas du pavé). \deleted{Donc} \deleted{dans ce} \unsure{a)} Si $\nu=0$, $u$ est injectif donc \uline{iso}, et \deleted{\unsure{l'image}}\inserted{les op.} de $V$ sur $E$ via $\rho$ ou via $\sigma$ se déterminent mutuellement. Or l'application $\varphi$ commute bien
